\section{模型的建立与求解方法}\label{sec:model}

\subsection{控制方程与定解条件}

取横截面上半径 $r$ 处的薄圆环，水分按 Fick 定律扩散\cite{crank}，热量按 Fourier 定律传导\cite{carslaw}，由圆环的水分与热量收支得到径向控制方程
\begin{gather}
\pd{C}{t}=\frac1r\pd{}{r}\!\left(rD(C,T)\pd{C}{r}\right),\qquad 0<r<R_0 \label{eq:govC}\\
\rho(C)c_p(C)\pd{T}{t}=\frac1r\pd{}{r}\!\left(rk(C)\pd{T}{r}\right) \label{eq:govT}
\end{gather}
方程取散度形式：扩散系数、导热系数随位置变化时，相邻圆环之间的水分通量 $-D\,\partial C/\partial r$ 与热通量 $-k\,\partial T/\partial r$ 连续，水分和热量整体守恒。各问的物性为：问题一 $\rho=820$、$c_p=2600$、$k=0.36$（单位见符号说明），$D=7\times10^{-9}\mathrm{e}^{-0.89/C}$；问题二、三按附录 3，$\rho=650+128C$，$c_p=1450+2736\,C/(C+1)$，$k=0.21+0.38\,C/(C+1)$，$D=2.4\times10^{-3}\mathrm{e}^{-0.45/C}\mathrm{e}^{-3850/T}$；问题四按附录 4，$\rho=760+90C$，$c_p=1850+2150\,C/(C+1)$，$k=0.12+0.20\,C/(C+1)$，$D=4.2\times10^{-4}\mathrm{e}^{-0.30/C}\mathrm{e}^{-3850/T}$。

初始时药材内部均匀，中心满足对称条件，表面与烘房空气之间按对流换热和对流传质交换：
\begin{gather}
T(r,0)=\SI{301.15}{K},\qquad C(r,0)=C_0=2.55 \label{eq:init}\\
\pd{C}{r}\bigg|_{r=0}=\pd{T}{r}\bigg|_{r=0}=0 \label{eq:center}\\
-D\pd{C}{r}\bigg|_{r=R_0}=h_m\big[C_s-\Cenv(t)\big],\qquad -k\pd{T}{r}\bigg|_{r=R_0}=h\big[T_s-\Tair(t)\big] \label{eq:bc}
\end{gather}
烘房温度 $\Tair(t)$ 与含水浓度 $\Cenv(t)$ 在 $0$--$\SI{4}{h}$ 内按附件 1（每 $\SI{60}{s}$ 一个数据）线性插值；$\SI{4}{h}$ 后取最后 $\SI{1}{h}$ 共 61 个数据的平均值 $\SI{49.9989}{\celsius}$ 和 $0.049988$（图~\ref{fig:environment}）。

\begin{figure}[htbp]\centering
\includegraphics[width=\textwidth]{fig_environment.pdf}
\caption{烘房的湿空气状态。(a) 烘房温度与湿球温度；(b) 含水浓度与相对湿度。$\SI{4}{h}$ 后（阴影区）取最后 $\SI{1}{h}$ 数据的平均值；湿球温度与相对湿度按第~\ref{sec:evap}~节的湿空气关系由含水浓度折算。}
\label{fig:environment}
\end{figure}

将式~\eqref{eq:govC} 乘以 $2r/R_0^2$ 在截面上积分，并代入式~\eqref{eq:bc}，得到截面平均含水率 $\Cbar=\frac{2}{R_0^2}\int_0^{R_0}Cr\,\mathrm{d}r$ 的变化率
\begin{equation}
-\frac{\mathrm{d}\Cbar}{\mathrm{d}t}=\frac{2h_m}{R_0}\big[C_s-\Cenv(t)\big]
\label{eq:rate}
\end{equation}
它就是药材的干燥速率，只由表面含水率与烘房含水浓度之差决定。

\subsection{问题四的参考坐标}\label{sec:ref}

问题四的半径 $R(t)$ 由附件 2 给出（每 $\SI{0.5}{h}$ 一个数据，$\SI{72}{h}$ 内由 $\SI{2}{cm}$ 降到 $\SI{1.198}{cm}$），线性插值。药材均匀收缩时，干物质随骨架一起移动，干物质密度 $\rhos=\rhos^0(R_0/R)^2$ 只随时间变化。取随干物质移动的参考坐标 $x=r/R(t)\in[0,1]$：干基含水率是单位干物质的含水量，沿干物质轨迹的变化率就是 $x$ 固定时的偏导数，水分守恒式减去干物质守恒式后对流项消去\cite{crank1984}，得到
\begin{gather}
\pd{C}{t}=\frac{1}{R^2}\frac1x\pd{}{x}\!\left(xD\pd{C}{x}\right),\qquad
\rho c_p\pd{T}{t}=\frac{1}{R^2}\frac1x\pd{}{x}\!\left(xk\pd{T}{x}\right),\qquad 0<x<1 \label{eq:ref}\\
-\frac{D}{R}\pd{C}{x}\bigg|_{x=1}=h_m\big[C_s-\Cenv\big],\qquad
-\frac{k}{R}\pd{T}{x}\bigg|_{x=1}=h\big[T_s-\Tair\big] \label{eq:refbc}
\end{gather}
计算域固定为 $[0,1]$，收缩只体现为扩散项的系数 $1/R^2$ 和边界项的系数 $1/R$（图~\ref{fig:geometry}(b)），式~\eqref{eq:rate} 中的 $R_0$ 相应换为 $R(t)$。题目要求的固定位置 $r_j$ 在 $t$ 时刻对应 $x_j=r_j/R(t)$；当 $r_j>R(t)$ 时该位置已在药材之外，结果留空。

\begin{figure}[htbp]\centering
\includegraphics[width=\textwidth]{fig_geometry.pdf}
\caption{物理模型示意。(a) 横截面上的热量、水分交换与中心对称条件，同心环的深浅示意内湿外干的含水率分布；(b) 问题四中物理域 $[0,R(t)]$ 随附件 2 收缩，映射到固定的参考域 $[0,1]$ 上计算。}
\label{fig:geometry}
\end{figure}

\subsection{数值求解方法}\label{sec:numeric}

\parahead{空间离散} 采用节点型有限体积法\cite{patankar}。把 $[0,R_0]$（问题四为 $[0,1]$）等分为 $N$ 个区间，节点 $r_i=i\Delta r$，中心和表面都是节点，题目要求的输出位置都落在节点上。对式~\eqref{eq:govC} 在节点 $i$ 的控制体 $[r_{i-1/2},r_{i+1/2}]$ 上积分得
\begin{equation}
V_i\frac{\mathrm{d}C_i}{\mathrm{d}t}=r_{i+\frac12}D_{i+\frac12}\frac{C_{i+1}-C_i}{\Delta r}-r_{i-\frac12}D_{i-\frac12}\frac{C_i-C_{i-1}}{\Delta r},\qquad V_i=\int_{r_{i-1/2}}^{r_{i+1/2}}r\,\mathrm{d}r
\label{eq:fvm}
\end{equation}
中心和表面节点只有半个控制体，表面节点外侧的通量由式~\eqref{eq:bc} 给出，温度方程的离散相同。扩散系数随含水率变化剧烈（含水率由 2.55 降到 0.15 时减小约 17 倍），界面扩散系数取两侧含水率之间的积分平均 $D_{i+1/2}=\int_0^1D\big((1-\xi)C_i+\xi C_{i+1},\bar T_{i+1/2}\big)\mathrm{d}\xi$（8 点 Gauss 求积），在陡峭的含水率前沿处比算术平均、调和平均更稳健\cite{maddix}；导热系数变化平缓，取调和平均\cite{patankar}。离散后得到 $2(N+1)$ 维常微分方程组。

\parahead{时间积分} 热量与水分的时间尺度相差一到两个数量级，扩散系数又随时间大幅变化，方程组刚性很强。采用变阶、变步长的 BDF 方法\cite{shampine}（SciPy 的 \texttt{solve\_ivp}\cite{scipy}），相对容差 $10^{-8}$，含水率与温度的绝对容差分别为 $10^{-11}$ 和 $\SI{e-8}{K}$；前 $\SI{4}{h}$ 的最大步长取 $\SI{60}{s}$，与附件 1 的数据间隔一致，之后放宽到 $\SI{600}{s}$，并在 $\SI{4}{h}$ 处重新起步，避免跨越输入数据的折点。网格取 $N=1600$（问题一）和 $N=800$（问题二至四），即 $\Delta r=\SI{0.0125}{mm}$ 和 $\SI{0.025}{mm}$；输出时刻的数值由 BDF 的稠密输出给出。

\parahead{烘干时长} 记全域最大含水率 $\Cmax(t)=\max_iC_i(t)$。烘干时长 $\tstar$ 定义为 $\Cmax$ 首次降到 $0.15$ 的时刻，在积分过程中以 $\Cmax(t)-0.15=0$ 为事件函数精确定位。result3、result4 每隔 $\SI{60}{s}$ 输出一行，直到首个最大含水率低于 $0.15$ 的采样时刻。
